%% kapa: exam=1 round=25 year=2014 period=1 form=A subject=economics subject_ko=경제원론 seq=11 number=51
%% source: Q-Net 감정평가사 기출문제 (cst00309) artlSeq=5101207 file=extracted/1008/5101207/2014년 제25회 감정평가사 1차시험 기출문제/1교시(A형).hwp
%% conversion: kordoc 4.24.51 -> tools/gen1.py (auto); review=proofread-sonnet; flags=none
\begin{question}{51}
A, B 두 나라에 각각 다섯 사람씩 살고 있다고 한다. A국과 B국에 사는 사람들의 소득은 각각 (1, 1, 2, 2, 4), (1, 2, 2, 2, 3)이라고 한다. 다음 설명 중 옳은 것은?
\begin{choices}
  \item 지니계수상으로 A국이 더 평등하며, 십분위분배율로 보면 B국이 더 평등하다.
  \item 지니계수상으로 B국이 더 평등하며, 십분위분배율로 보면 A국이 더 평등하다.
  \item 지니계수상으로 B국이 더 평등하며, 십분위분배율은 현재의 정보로는 계산할 수 없다.
  \item 지니계수상으로나 십분위분배율로나 A국이 더 평등하다.
  \item 지니계수상으로나 십분위분배율로나 B국이 더 평등하다.
\end{choices}
\end{question}
\answer{5}
